% TEMPLATE-MILC-v3.tex - plantilla maestra UNICA de las guias MILC v3.
%
% La usan las cuatro familias de guias, cada una con su driver:
%   · Grados 6-11          -> scripts/build-guias-g11.py
%   · Bebras               -> scripts/build-guias-bebras.py
%   · Semillero            -> scripts/build-guias-semillero.py
%   · Territorio interior  -> scripts/build-guias-territorio-interior.py
%
% Antes habia tres copias casi identicas (~400 lineas iguales, 6 distintas):
% cada mejora habia que aplicarla tres veces, y en la practica no se hacia
% (el motor de recursos vivia solo en dos de ellas). Lo que varia entre
% familias esta parametrizado en 6 placeholders que cada driver llena
% (se nombran aqui SIN su sintaxis real: el builder reemplaza por texto plano,
% tambien dentro de los comentarios, y un valor multilinea rompe el preambulo):
%
%   HEADER_IZQ        encabezado izquierdo de pagina
%   HEADER_DER        encabezado derecho de pagina
%   PORTADA_ETIQUETA  rotulo sobre el numero grande (GRADO/MOMENTO/MODULO)
%   PORTADA_SERIE     linea de serie de la portada
%   PORTADA_ID        identificador de la guia en la portada
%   PORTADA_CIRCULO   el circulito de grado (°), o vacio si no aplica
%   PDF_TITULO        titulo en texto plano (sin \textbf ni comillas TeX) para
%                     los metadatos del PDF (/Title); lo produce lib_xelatex.texto_pdf
%
% Composicion (auditoria 2026-09-03): las bandas de estacion piden espacio
% (\Needspace*) y prohiben el salto justo despues (\nopagebreak), asi no quedan
% huerfanas al pie; las cajas largas (softbox, drawbox, infoband, puentebox) son
% `breakable`, asi una caja de media pagina ya no arrastra todo a la siguiente
% dejando la anterior al 40 %. Todo texto sobre mostaza o verde va en milcNegro
% (blanco sobre #E5B400 da 1,93:1 y sobre #58A923 2,95:1; ninguno pasa WCAG AA).
% El resto de claves las documenta content/guias/_SCHEMA.md. El script valida
% que no queden placeholders sin reemplazar antes de escribir el archivo final.

% Etiquetado PDF (latex-lab): /Lang, estructura y texto alternativo de las
% figuras, para lectores de pantalla. Debe ir ANTES de \documentclass.
\DocumentMetadata{lang=es-CO,tagging=on}
\documentclass[11pt,letterpaper]{article}
% headheight=24pt: la cabecera derecha lleva dos lineas (identidad de la guia
% y estacion actual). headsep se acorta en la misma medida para que el bloque
% de texto quede exactamente donde estaba con headheight=12pt.
\usepackage[margin=2.0cm,headheight=24pt,headsep=13pt]{geometry}
\usepackage[table]{xcolor}
\usepackage{fontspec}
\usepackage[spanish]{babel}
\usepackage{tikz}
% Librerías TikZ para los diagramas de las guías (bloque `recursos:`):
%  · babel  → babel en español vuelve activos caracteres como «>», y TikZ los usa
%             en sus opciones (>=stealth, ->). Sin esta librería, cualquier
%             diagrama con flechas revienta con «\language@active@arg> has an
%             extra }». Debe ir ANTES de que se dibuje nada.
%  · positioning        → sintaxis «below left=2pt and 3pt».
%  · decorations.pathreplacing → llaves ({brace}) para señalar rangos.
\usetikzlibrary{babel,positioning,decorations.pathreplacing}
\usepackage{graphicx}
% Circuitos: el trabajo de electrónica se hace hoy con micro:bit y sus sensores
% integrados (MakeCode, sin cableado), así que no se necesitan esquemáticos.
% Si algún día entran montajes con protoboard, basta descomentar esta línea:
% los diagramas .tex/.tikz declarados en `recursos:` ya se incrustan solos.
% \usepackage{circuitikz}
\usepackage[most]{tcolorbox}
\usepackage{tabularx}
\usepackage{array}
\usepackage{fancyhdr}
\usepackage{titlesec}
\usepackage[document]{ragged2e}  % bandera a la derecha en todo el cuerpo
\usepackage{enumitem}
\usepackage{microtype}
\usepackage{etoolbox}
\usepackage{needspace}
% hyperref al final del preambulo: metadatos del PDF (titulo, autor, idioma,
% familia), marcadores por estacion (\pdfbookmark) y enlaces sin recuadro.
\usepackage[hidelinks,
  pdftitle={Encabezado, pie de página y números — terminar el documento como un profesional},
  pdfauthor={PhD. Álvaro Cárdenas Orozco},
  pdfsubject={Tecnología e Informática · Grado 6 · Guía 26 · MILC v3},
  pdflang={es-CO},
  bookmarksopen=true,bookmarksnumbered=false]{hyperref}

% Helvetica Neue no trae la flecha U+2192, asi que cada «→» escrito literalmente
% en el YAML se compilaba a NADA, en silencio (462 flechas en 61 guias: rutas del
% tipo «paleta → tipografia → lockup» salian rotas). Mapeamos el caracter a su
% version matematica, que si tiene glifo. Asi el autor puede seguir escribiendo
% «→» comodamente en el YAML.
\usepackage{newunicodechar}
\newunicodechar{→}{\ensuremath{\rightarrow}}
% Mismo problema con los simbolos de marca y vinetas: el «✓» de «una palomita ✓»
% salia como cuadrito vacio en el PDF de todas las guias que lo usaban.
\usepackage{pifont}
\newunicodechar{✓}{\ding{51}}
\newunicodechar{✗}{\ding{55}}
\newunicodechar{✔}{\ding{52}}
\newunicodechar{•}{\textbullet}
\newunicodechar{≥}{\ensuremath{\geq}}
\newunicodechar{≤}{\ensuremath{\leq}}

\setmainfont{Helvetica Neue}
\setsansfont{Helvetica Neue}
\setlength{\parindent}{0pt}
\setlength{\parskip}{6pt}
% Sin particion de palabras: en 6.º-7.º una palabra cortada por guion al final
% de linea («Comuni-cacion») frena la lectura mucho mas que una linea corta.
\hyphenpenalty=10000
\exhyphenpenalty=10000
\pretolerance=2000
\tolerance=3000
\setlength{\emergencystretch}{3em}
\linespread{1.10}
\renewcommand{\arraystretch}{1.25}
\setlist{leftmargin=5mm,itemsep=1mm,topsep=1mm}
\newcolumntype{Y}{>{\RaggedRight\arraybackslash}X}

\definecolor{milcMagenta}{HTML}{E6007E}
\definecolor{milcVino}{HTML}{5A0038}
\definecolor{milcTurquesa}{HTML}{0093A5}
\definecolor{milcVerde}{HTML}{58A923}
\definecolor{milcMostaza}{HTML}{E5B400}
\definecolor{milcOcre}{HTML}{B86600}
\definecolor{milcNegro}{HTML}{241816}
\definecolor{milcGris}{HTML}{F7F7F4}
\definecolor{milcLinea}{HTML}{E8E2DD}
\definecolor{milcGradePrimary}{HTML}{0066FF}
\definecolor{milcGradeDark}{HTML}{003D99}
\definecolor{milcGradeSoft}{HTML}{D6E8FF}
\definecolor{milcGradeAccent}{HTML}{CFE7FF}
\definecolor{milcGradeText}{HTML}{FFFFFF}
\color{milcNegro}

\pagestyle{fancy}
\fancyhf{}
% Cabecera de dos lineas: identidad de la guia arriba y, a la derecha y en
% gris, la estacion en curso (\markright desde \milcestacion). Antes de la
% primera estacion la segunda linea va vacia; el \strut mantiene la altura
% para que la linea de arriba no baile entre paginas.
\lhead{\small Tecnología e Informática · Grado 6\\[-1pt]{\footnotesize\strut}}
\rhead{\small Guía 26 · MILC v3\\[-1pt]{\footnotesize\strut\color{milcNegro!65}\rightmark}}
\cfoot{\small\thepage}
\renewcommand{\headrulewidth}{0pt}

% Sobre mostaza (#E5B400) y verde (#58A923) el texto va en milcNegro; sobre el
% resto de la paleta, en blanco. Se usa dentro de fonttitle/fontupper, donde un
% \color posterior manda sobre coltitle.
\newcommand{\milctintasobre}[1]{%
  \ifboolexpr{test{\ifstrequal{#1}{milcMostaza}} or test{\ifstrequal{#1}{milcVerde}}}%
    {\color{milcNegro}}{\color{white}}}

\newtcolorbox{titlebox}[1]{
  enhanced,colback=#1,colframe=#1,arc=7mm,boxrule=0pt,
  left=7mm,right=7mm,top=3mm,bottom=3mm,
  before skip=8pt,after skip=8pt,
  fontupper=\sffamily\bfseries\Large\color{white}
}

\newtcolorbox{softbox}[3]{
  enhanced,breakable,pad at break=3mm,
  colback=#2,colframe=#1,borderline west={4pt}{0pt}{#1},
  arc=2mm,boxrule=.4pt,left=4mm,right=4mm,top=3mm,bottom=3mm,
  before skip=6pt,after skip=8pt,title={#3},coltitle=white,
  colbacktitle=#1,fonttitle=\sffamily\bfseries\milctintasobre{#1},fontupper=\RaggedRight,
  attach boxed title to top left={xshift=3mm,yshift=-2mm},
  boxed title style={arc=2mm, boxrule=0pt}
}

\newtcolorbox{drawbox}[2]{
  enhanced,breakable,pad at break=4mm,
  colback=white,colframe=#1,arc=4mm,boxrule=1pt,
  left=5mm,right=5mm,top=4mm,bottom=4mm,before skip=8pt,after skip=8pt,
  title={#2},coltitle=white,colbacktitle=#1,fonttitle=\sffamily\bfseries\milctintasobre{#1},
  fontupper=\RaggedRight,
  attach boxed title to top left={xshift=4mm,yshift=-2mm},
  boxed title style={arc=3mm, boxrule=0pt}
}

\newtcolorbox{infoband}[1]{
  enhanced,breakable,pad at break=2mm,colback=#1!8,colframe=#1!50,arc=2mm,boxrule=.5pt,
  left=4mm,right=4mm,top=2mm,bottom=2mm,before skip=4pt,after skip=4pt,
  fontupper=\small\RaggedRight
}

% Puente narrativo entre secciones (contrato editorial Conéctate).
% Da continuidad: "Acabas de X. Ahora vas a Y porque Z."
% Visualmente: banda izquierda en color del grado, texto en itálica pequeña.
\newtcolorbox{puentebox}{
  enhanced,breakable,pad at break=2mm,colback=milcGradeSoft,colframe=milcGradeDark,
  borderline west={3pt}{0pt}{milcGradeDark},
  arc=0mm,boxrule=0pt,left=5mm,right=4mm,top=2.5mm,bottom=2.5mm,
  before skip=4pt,after skip=8pt,
  fontupper=\itshape\small\color{milcNegro}\RaggedRight
}
% La flecha va en modo matematico a proposito: Helvetica Neue NO tiene el glifo
% U+2192, asi que \textrightarrow se compilaba a NADA («Missing character: There
% is no → in font Helvetica Neue») y el puente salia sin flecha. En modo
% matematico la flecha viene de la fuente matematica y si se dibuja.
\newcommand{\puente}[1]{%
  \begin{puentebox}%
    \textcolor{milcGradeDark}{$\rightarrow$}\hspace{2mm}#1%
  \end{puentebox}%
}

\newcommand{\linead}[1]{\noindent\color{milcLinea}\rule{#1}{0.5pt}\color{milcNegro}}
\newcommand{\respuesta}[1]{\par\vspace{2mm}\foreach \n in {1,...,#1}{\noindent\color{milcLinea}\rule{\linewidth}{0.5pt}\par\vspace{5mm}}\color{milcNegro}}
\newcommand{\checkbox}{\textcolor{milcGradeDark}{\fbox{\phantom{x}}}\hspace{5pt}}

% ─── Listas aireadas para las guías socioemocionales ────────────────────
% Componer las verificaciones como «\checkbox texto\\» deja el texto que salta
% de línea debajo del cuadrito y apelmaza el bloque. Estos entornos alinean la
% continuación y airean los ítems. Los usa Territorio Interior; cualquier
% familia puede hacerlo.
\newlist{checklista}{itemize}{1}
\setlist[checklista]{
  label=\textcolor{milcGradeDark}{\fbox{\phantom{x}}},
  leftmargin=9mm, labelsep=3.2mm, itemsep=2.4mm,
  topsep=2.5mm, parsep=0pt, partopsep=0pt, before=\RaggedRight
}
\newlist{pasoslista}{enumerate}{1}
\setlist[pasoslista]{
  label=\textcolor{milcGradeDark}{\sffamily\bfseries\arabic*.},
  leftmargin=8mm, labelsep=3mm, itemsep=2.8mm,
  topsep=1mm, parsep=0pt, partopsep=0pt, before=\RaggedRight
}

\newcommand{\phasecard}[4]{
\begin{tcolorbox}[enhanced,colback=#3,colframe=#2,arc=3mm,boxrule=.5pt,height=3.05cm,valign=center,halign=center,left=1.5mm,right=1.5mm,top=2mm,bottom=2mm]
{\sffamily\bfseries\fontsize{22}{24}\selectfont\textcolor{#2}{#1}}\par
{\sffamily\bfseries\small #4}
\end{tcolorbox}
}

% ── Piezas del rediseno 2026 (legibilidad 6.º-11.º) ───────────────────
% Banda de estacion: reemplaza el titlebox macizo. Titulo a la izquierda,
% dato practico (tiempo/modalidad) a la derecha, en una sola linea.
% Nunca huerfana: pide 9 lineas libres (o salta de pagina rellenando con
% \vfill) y prohibe el salto justo despues de la banda. Nueve y no seis:
% la caja partible que sigue necesita ~80pt (titulo adosado, relleno y dos
% lineas) para arrancar en la misma pagina; con menos, tcolorbox la manda
% entera a la siguiente y la banda se queda sola. El penalty 10000 va
% en `after`, ANTES del espacio posterior: un `after skip` (aunque sea 0pt)
% es un \vskip tras la caja, y ese glue es un punto de corte legal que un
% \nopagebreak escrito despues de \end{tcolorbox} ya no cierra. Cada banda
% deja marcador PDF y actualiza la cabecera derecha.
\newcounter{milcestacion}
\newcommand{\milcestacion}[2]{%
\Needspace*{9\baselineskip}%
\stepcounter{milcestacion}%
\pdfbookmark[1]{#2}{milcestacion\themilcestacion}%
\markright{#2}%
\begin{tcolorbox}[enhanced,colback=#1,colframe=#1,arc=2mm,boxrule=0pt,
  left=5mm,right=5mm,top=2.6mm,bottom=2.6mm,before skip=11pt,
  after={\par\nopagebreak\vskip7pt}]
{\sffamily\bfseries\fontsize{15}{18}\selectfont\milctintasobre{#1}#2}
\end{tcolorbox}}

% Tarjeta de paso numerada: sustituye las tablas «Pilar / Aplicacion», que
% eran formato de consulta docente, no de estudiante. El numero va en un
% circulo sobre el borde para que la secuencia se lea de un vistazo.
\newcommand{\milcpaso}[3]{%
\Needspace{4\baselineskip}%
\begin{tcolorbox}[enhanced,colback=white,colframe=milcLinea,arc=1.5mm,boxrule=.9pt,
  left=14mm,right=4mm,top=3mm,bottom=3mm,before skip=4pt,after skip=4pt,
  fontupper=\RaggedRight,
  overlay={\node[anchor=north west,circle,fill=#2,text=white,inner sep=0pt,
    minimum size=7.4mm,font=\sffamily\bfseries\small]
    at ([xshift=3.6mm,yshift=-3.2mm]frame.north west) {#1};}]
#3
\end{tcolorbox}}

% Casilla marcable de tamano util con lapiz (la anterior era \fbox{\phantom{x}},
% de alto variable segun el texto que la seguia).
\newcommand{\milcmarca}{\framebox[3.8mm]{\rule{0pt}{3.4mm}}}

% Fila marcable de lista suelta (checks de la estacion 1).
\newcommand{\milccheckrow}[1]{%
\par\vspace{1.8mm}\noindent
\makebox[6.4mm][l]{\raisebox{-.55mm}{\textcolor{milcNegro}{\milcmarca}}}%
\parbox[t]{\dimexpr\linewidth-6.4mm\relax}{\RaggedRight #1}\par}

% Celda marcable dentro de la rubrica (tabla de autoevaluacion). Misma
% sangria francesa, pero pensada para vivir dentro de una columna X.
\newcommand{\milcopcion}[1]{%
\makebox[5.2mm][l]{\raisebox{-.55mm}{\textcolor{milcNegro}{\milcmarca}}}%
\parbox[t]{\dimexpr\linewidth-5.2mm\relax}{\RaggedRight #1}}

% ── Recursos visuales de la guía ──────────────────────────────────────
% Declarados en el bloque `recursos:` del YAML y emitidos por el builder.
% \guiaFigura   → imagen rasterizada o PDF (foto, carta celeste, montaje).
% \guiaDiagrama → fuente TikZ/circuitikz (.tex) incrustada como vector:
%                 es la vía para circuitos y esquemáticos editables.
% [#1] = texto alternativo (alt) · #2 = ruta relativa al .tex · #3 = pie
% El alt llega al PDF etiquetado (/Alt) y lo lee el lector de pantalla.
\newcommand{\guiaFigura}[3][]{%
\begin{center}
\includegraphics[width=0.88\linewidth,height=0.38\textheight,keepaspectratio,alt={#1}]{#2}\\[1.5mm]
{\footnotesize\itshape\textcolor{milcNegro!75}{#3}}
\end{center}
}

\newcommand{\guiaDiagrama}[2]{%
\begin{center}
\input{#1}\\[1.5mm]
{\footnotesize\itshape\textcolor{milcNegro!75}{#2}}
\end{center}
}

\begin{document}
\thispagestyle{empty}

% ── Portada ───────────────────────────────────────────────────────────
% Antes: un rectangulo de color a pagina completa. Se veia bien en pantalla,
% pero en la fotocopiadora del colegio se comia el toner y salia gris sucio.
% Ahora la banda ocupa el tercio superior y el resto va en blanco.
\begin{tikzpicture}[remember picture,overlay]
  \path[fill=milcGradePrimary]
    (current page.north west) rectangle ([yshift=-10.6cm]current page.north east);

  \node[anchor=north west,text=milcGradeText] at ([xshift=2.0cm,yshift=-1.55cm]current page.north west)
    {\sffamily\bfseries\fontsize{12}{14}\selectfont GRADO};
  \node[anchor=north west,text=milcGradeText,opacity=.85] at ([xshift=2.0cm,yshift=-2.35cm]current page.north west)
    {\sffamily\bfseries\fontsize{8.5}{10}\selectfont SERIE GUÍAS · TECNOLOGÍA E INFORMÁTICA};
  \node[anchor=north east,text=milcGradeText,opacity=.85,align=right] at ([xshift=-2.0cm,yshift=-1.55cm]current page.north east)
    {\sffamily\bfseries\fontsize{9}{12}\selectfont I.E. Sor María Juliana\\Cartago, Valle del Cauca};

  \node[anchor=west,text=milcGradeText] (gradonum) at ([xshift=1.85cm,yshift=-7.1cm]current page.north west)
    {\sffamily\bfseries\fontsize{112}{112}\selectfont 6};
  % El circulito de grado (°) solo aplica a las guias de grado («11°»); en
  % Bebras y Semillero el numero grande es un momento/modulo, no un grado.
  % Se posiciona relativo al nodo `gradonum`, no en coordenadas absolutas.
  \draw[milcGradeText,line width=1.6pt]
    ([xshift=2.0mm,yshift=-4.5mm]gradonum.north east) circle (.15cm);

  \node[anchor=west,text=milcGradeText,text width=11.4cm,align=left] at ([xshift=1.5cm,yshift=0.5cm]gradonum.east)
    {
     {\sffamily\bfseries\fontsize{9}{11}\selectfont\MakeUppercase{Período 3 · Procesador de texto e Internet}}\\[3mm]
     {\sffamily\bfseries\fontsize{21}{25}\selectfont\setlength{\baselineskip}{27pt}Encabezado y\\[4pt]pie de página ---\\[4pt]terminar el\\[4pt]documento como\\[4pt]un profesional}\\[3.5mm]
     {\sffamily\bfseries\fontsize{15}{18}\selectfont Guía 26-6-TIC}
    };
\end{tikzpicture}

\vspace*{9.4cm}
\vfill

{\sffamily\bfseries\fontsize{8}{10}\selectfont\textcolor{milcGradeDark}{LO QUE TE LLEVAS HOY}}

\begin{tcolorbox}[enhanced,colback=white,colframe=milcNegro,arc=2mm,boxrule=1.6pt,
  left=5mm,right=5mm,top=4mm,bottom=4mm,before skip=3pt,after skip=12pt,
  fontupper=\RaggedRight\fontsize{11}{15.5}\selectfont]
Un \textbf{mini-informe de 2 páginas} sobre cualquier tema (puede ser tu colegio, tu barrio, una materia favorita), formateado con: \textbf{(a) encabezado} arriba (con tu nombre, grado y materia); \textbf{(b) pie de página} con la fecha y la frase \emph{``Plataforma Conéctate''}; \textbf{(c) número de página} en la esquina inferior derecha (1, 2); \textbf{(d) portada de 1 página adicional} con título grande, tu nombre, fecha, grado; (e) el cuerpo del informe en las 2 páginas siguientes. Más \textbf{en el cuaderno}: boceto del layout con los 4 elementos etiquetados.
\end{tcolorbox}

\begin{tcolorbox}[enhanced,colback=milcGris,colframe=milcGris,arc=2mm,boxrule=0pt,
  left=5mm,right=5mm,top=3.5mm,bottom=3.5mm,after skip=12pt]
{\sffamily\bfseries\fontsize{8}{10}\selectfont\textcolor{milcNegro!55}{ESTA GUÍA ES DE}}\\[3mm]
\begin{tabularx}{\linewidth}{X p{3.2cm} p{3.2cm}}
\linead{\linewidth} & \linead{3.0cm} & \linead{3.0cm}\\[-1mm]
{\fontsize{8.5}{10}\selectfont\textcolor{milcNegro!55}{Nombre completo}} &
{\fontsize{8.5}{10}\selectfont\textcolor{milcNegro!55}{Grupo}} &
{\fontsize{8.5}{10}\selectfont\textcolor{milcNegro!55}{Fecha}}\\
\end{tabularx}
\end{tcolorbox}

\vfill

\noindent\textcolor{milcLinea}{\rule{\linewidth}{1pt}}\\[2mm]
{\sffamily\bfseries\fontsize{9.5}{12}\selectfont Autor: Dr. Álvaro Cárdenas Orozco}\\[1mm]
{\sffamily\fontsize{8.5}{11}\selectfont\textcolor{milcNegro!55}{Metodología MILC · apertura ancestral · pensamiento computacional · triángulo Dussel–estoicismo–Floridi}}

\newpage

\section*{Apertura: Encabezado, pie de página y números --- terminar el documento como un profesional}

\begin{softbox}{milcGradeDark}{milcGradeSoft}{Saber ancestral}
\textbf{Cuando doña Mercedes hacía un cuaderno especial, le ponía marco al trabajo.} En la escuela rural de la vereda La Plata de Cartago, doña Mercedes la maestra exigía que los cuadernos del campo tuvieran \textbf{el nombre arriba} (\emph{``Cuaderno de Lengua, Pedro Martínez, grado 5°''}) y \textbf{la fecha abajo} (\emph{``Año 1985''}). En cada hoja, los niños debían escribir \textbf{su nombre arriba} (\emph{el encabezado}) y \textbf{el número de la hoja abajo} (\emph{el pie}). Eso no era manía: era \emph{darle marco al trabajo}. Si un cuaderno se caía o un niño olvidaba el suyo, todos podían identificarlo en segundos. Si una hoja se desprendía, se podía colocar de nuevo gracias al número. \textbf{Doña Mercedes enseñaba con eso un saber profundo: todo trabajo serio tiene marco}. Hoy en digital los procesadores te dan herramientas para hacer ese marco: encabezado, pie de página y números. Y se aplican a CUALQUIER documento de varias páginas que vayas a entregar en tu vida.
\end{softbox}

\begin{softbox}{milcTurquesa}{milcGris}{Saber contemporáneo}
Cuando un documento tiene \textbf{más de 1 página}, necesita \textbf{3 elementos de marco}. \textbf{(1) Encabezado} (\emph{header}): la franja arriba de cada página, donde van tu nombre, materia, fecha o título del documento. Se ve en TODAS las páginas. \textbf{(2) Pie de página} (\emph{footer}): la franja abajo de cada página, donde va información secundaria (fecha de impresión, nombre del colegio, frase de cierre). \textbf{(3) Número de página}: el contador que dice qué página estás viendo (1, 2, 3...). Se inserta en el encabezado o el pie. \textbf{Estos 3 elementos te ahorran problemas:} si entregas 5 hojas y se desordenan, el número resuelve. Si el profe pierde tu trabajo entre 30 entregas, el encabezado con tu nombre lo identifica. Es como ponerle marco a un cuadro: \emph{el contenido sigue siendo el cuadro, pero el marco le da presencia}.
\end{softbox}

\begin{softbox}{milcMostaza}{milcGris}{Pregunta-puente}
\textit{Tu profe acumula \textbf{30 trabajos impresos} en su mesa. Las hojas se mezclan en el escritorio. ¿\textbf{Cómo se asegura el profe} de identificar el tuyo? ¿Y si tu trabajo tenía 4 hojas y se desordenaron, cómo sabe qué hoja iba primero?}
\end{softbox}

\begin{softbox}{milcVerde}{milcGris}{Saber y saber hacer}
Al terminar la clase: (1) podrás \textbf{identificar} los 3 elementos de marco (encabezado, pie, número); (2) sabrás \textbf{aplicar} los 3 en Word/Docs; (3) podrás \textbf{distinguir} entre portada y página de contenido; (4) habrás \textbf{creado} un mini-informe de 2 páginas con todos los elementos.
\end{softbox}



\small
\begin{tabularx}{\textwidth}{>{\bfseries}p{3.0cm}Y}
\rowcolor{milcGradePrimary!12}Autor & Dr. Álvaro Cárdenas Orozco\\
DBA & Produce contenidos digitales y valida información de fuentes web (MEN, T\&I 6°)\\
Referentes & Saber ancestral de la maestra rural · Lectura crítica · Soberanía informacional\\
Producto final & Un \textbf{mini-informe de 2 páginas} sobre cualquier tema (puede ser tu colegio, tu barrio, una materia favorita), formateado con: \textbf{(a) encabezado} arriba (con tu nombre, grado y materia); \textbf{(b) pie de página} con la fecha y la frase \emph{``Plataforma Conéctate''}; \textbf{(c) número de página} en la esquina inferior derecha (1, 2); \textbf{(d) portada de 1 página adicional} con título grande, tu nombre, fecha, grado; (e) el cuerpo del informe en las 2 páginas siguientes. Más \textbf{en el cuaderno}: boceto del layout con los 4 elementos etiquetados.\\
\end{tabularx}
\normalsize

\puente{Hoy haces 4 cosas: descubres dónde están encabezado y pie en Word, aprendes a insertarlos, agregas números de página, y armas tu informe de 2 páginas.}

\milcestacion{milcGradeDark}{Ruta MILC: así vamos a aprender}
\begin{tabularx}{\textwidth}{>{\centering\arraybackslash}X>{\centering\arraybackslash}X>{\centering\arraybackslash}X>{\centering\arraybackslash}X}
\phasecard{1}{milcVerde}{milcGris}{Escucha\\[-1mm]\small Observo, escucho y pregunto.} &
\phasecard{2}{milcTurquesa}{milcGris}{Entender\\[-1mm]\small Sistematizo y ordeno.} &
\phasecard{3}{milcMagenta}{milcGris}{Hacer\\[-1mm]\small Construyo y aplico.} &
\phasecard{4}{milcVino}{milcGris}{Cierre\\[-1mm]\small Evalúo y reflexiono.}
\end{tabularx}


\puente{Empezamos con un documento bien formateado para identificar sus 3 elementos de marco.}

\milcestacion{milcVerde}{Estación 1 de 3 · Escucha}

\begin{softbox}{milcVerde}{milcGris}{Escena para escuchar}
\textbf{Actividad 1 · IDENTIFICA --- Encuentra los 3 elementos de marco} (10 min · individual). El profe te muestra un documento profesional (puede ser una hoja impresa, una circular del colegio o una imagen en pantalla). Tu tarea: en el cuaderno, dibuja una hoja A4 (rectángulo) y \textbf{marca con flechas} dónde aparecen estos 3 elementos: \emph{encabezado}, \emph{pie de página}, \emph{número de página}. Si encuentras los 3, ponte \texttt{¡EXCELENTE!} al lado. Si solo encuentras 2, escribe cuál falta. \emph{Si no tienes documento físico hoy}, dibuja basado en cómo crees que se ve un documento profesional.
\end{softbox}

{\sffamily\bfseries\fontsize{8}{10}\selectfont\textcolor{milcNegro!55}{MARCA CUANDO LO HAYAS HECHO}}
\milccheckrow{Dibujé una hoja A4 en el cuaderno.}
\milccheckrow{Marqué con flechas los 3 elementos (encabezado, pie, número).}
\milccheckrow{Identifiqué cuántos elementos vi en el documento del profe.}

\begin{infoband}{milcVerde}


\textbf{Lo que va al cuaderno (Actividad 1 · IDENTIFICA).} Título: «Actividad 1 --- Los 3 elementos de marco». \textbf{Formato:} dibujo de hoja A4 con 3 flechas y nombres. \textbf{Extensión:} 1/3 de página. \textbf{Sabes que terminaste cuando} tienes los 3 elementos ubicados en tu dibujo con su nombre.
\end{infoband}

\puente{Después aprendes los 3 caminos en Word para encabezado, pie y números.}

\milcestacion{milcTurquesa}{Estación 2 de 3 · Entender}

\begin{softbox}{milcTurquesa}{milcGris}{Lo que vas a entender}
Ya viste dónde van. Ahora aprendes \textbf{cómo insertarlos} en Word y cuándo usar cada uno.
\end{softbox}

\milcpaso{1}{milcTurquesa}{\textbf{Encabezado --- la franja de arriba.} En Word: \texttt{Insertar → Encabezado → escoge un estilo simple} (\emph{En blanco} es lo más limpio). Aparece un área en gris en la parte superior de la página: ahí escribes. \textbf{Qué poner:} tu nombre, grado, materia, título del trabajo. Ejemplo típico: \emph{``María Cárdenas · Grado 6B · Tecnología''}. \emph{Importante:} lo que escribes aparece en TODAS las páginas del documento. Por eso es genial: lo configuras una vez y se aplica a todo.}
\milcpaso{2}{milcTurquesa}{\textbf{Pie de página --- la franja de abajo.} En Word: \texttt{Insertar → Pie de página → En blanco}. Aparece un área en gris en la parte inferior. \textbf{Qué poner:} fecha de entrega, nombre del colegio, frase corta, número de página. Ejemplo: \emph{``I.E. Sor María Juliana · Cartago · Mayo 2026''}. Igual que el encabezado: aparece en todas las páginas.}
\milcpaso{3}{milcTurquesa}{\textbf{Número de página.} En Word: \texttt{Insertar → Número de página → escoge posición} (más común: \emph{Parte inferior, esquina derecha}). El número se inserta automáticamente y se actualiza solo cuando agregas o quitas páginas. \textbf{Pista profesional:} si el documento tiene portada, normalmente no se numera la portada. La página 1 es la primera página de contenido.}
\milcpaso{4}{milcTurquesa}{\textbf{Diferencia entre portada y página de contenido.} La \textbf{portada} es la primera página, con título grande, nombre, fecha, grado. NO lleva texto del informe. Las \textbf{páginas de contenido} son las que siguen. Estructura típica: \emph{página 1 (portada) + páginas 2 y 3 (contenido) + página 4 (cierre o referencias)}. En documentos cortos (1-2 páginas), no es obligatorio tener portada separada. En informes formales, sí.}

\begin{drawbox}{milcTurquesa}{3 elementos de marco + portada}
\textbf{Encabezado:} arriba, en todas las páginas. Nombre + grado + materia. \\[1mm]
\textbf{Pie de página:} abajo, en todas las páginas. Fecha + colegio + frase. \\[1mm]
\textbf{Número de página:} esquina inferior derecha. Empieza en 1 (sin contar portada). \\[1mm]
\textbf{Portada:} primera página separada, sin encabezado/pie de contenido.
\end{drawbox}

\begin{softbox}{milcMostaza}{milcGris}{Errores comunes y su corrección}
\textbf{Errores frecuentes con encabezado y pie.} (1) \emph{Olvidar el nombre en el encabezado}: 30 trabajos en la mesa del profe = ninguno se sabe de quién es. (2) \emph{Numerar la portada como página 1}: técnicamente la portada es 0 o sin número. La página 1 es la primera de contenido. (3) \emph{Mismo texto en encabezado y pie}: redundante. Encabezado para identificación; pie para datos secundarios. (4) \emph{Encabezado en mayúsculas con tamaño grande}: queda gritando. Tamaño 10-11 en encabezado y pie es lo profesional. (5) \emph{Información distinta en encabezado de cada página}: tiene que ser lo mismo en todas (eso es lo que el encabezado hace bien).
\end{softbox}

\begin{infoband}{milcTurquesa}


\textbf{Lo que va al cuaderno (Actividad 2 · EXPLICA).} Título: «Actividad 2 --- Caminos para insertar encabezado, pie y números». \textbf{Formato:} 3 bloques con el camino en Word (Insertar → ... → ...) + qué poner en cada uno. \textbf{Extensión:} 3 bloques. \textbf{Sabes que terminaste cuando} podrías guiar a un compañero paso por paso a insertar los 3.
\end{infoband}

\puente{Con todo claro, creas tu mini-informe completo con portada y 2 páginas.}

\milcestacion{milcMagenta}{Estación 3 de 3 · Hacer}

\begin{softbox}{milcMagenta}{milcGris}{Cómo vas a construir tu producto}
\textbf{Actividad 3 · APLICA --- Tu mini-informe de 2 páginas con marco completo} (32 min · individual). Vas a crear un mini-informe profesional sobre un tema libre. Estructura: \emph{Portada (1 página) + Contenido (2 páginas)}. Con encabezado, pie de página y números.
\end{softbox}

\milcpaso{1}{milcMagenta}{\textbf{Paso 1 --- Escoge tu tema.} En el cuaderno escribe el tema y 3 puntos clave que vas a desarrollar. Tema sugerido: \emph{``Mi materia favorita''}, \emph{``Una persona admiro''}, \emph{``Algo que me gustaría aprender''}, lo que tú quieras. 3 puntos clave que vas a desarrollar en el cuerpo.}
\milcpaso{2}{milcMagenta}{\textbf{Paso 2 --- Portada (página 1).} Abre documento nuevo. En la primera página: título grande centrado tamaño 28 + tu nombre completo + grado + materia + fecha. Aplica formato bonito (negrita, color azul oscuro). La portada \emph{no lleva} encabezado, pie ni número de página. Al final de la portada, presiona \texttt{Ctrl + Enter} para saltar a una nueva página.}
\milcpaso{3}{milcMagenta}{\textbf{Paso 3 --- Encabezado para páginas 2 y 3.} En la página 2 (donde está el cursor), ve a \texttt{Insertar → Encabezado → En blanco}. Escribe: \emph{``[Tu nombre completo] · Grado 6[letra] · Tecnología''}. Tamaño 10, color gris oscuro, alineación derecha. Sal del modo encabezado haciendo doble clic en el cuerpo.}
\milcpaso{4}{milcMagenta}{\textbf{Paso 4 --- Pie de página + número.} Ve a \texttt{Insertar → Pie de página → En blanco}. Escribe a la izquierda: \emph{``I.E. Sor María Juliana --- Cartago --- Mayo 2026''}. Después ve a \texttt{Insertar → Número de página → Parte inferior, esquina derecha}. Aparece el número 1 en la página 2 y 2 en la página 3. Si la portada también recibió número, configura para que no se numere (es una opción avanzada; el profe te puede ayudar).}
\milcpaso{5}{milcMagenta}{\textbf{Paso 5 --- Cuerpo del informe.} En la página 2 escribe un título arriba: \emph{``[Tema]''} (tamaño 14, negrita, azul). Debajo escribe 3-4 párrafos con tus 3 puntos clave. Aplica formato de S3 (justificado, sangría, negrita en lo importante). En la página 3 termina los párrafos (o agrega una mini-conclusión, una lista, o una imagen). Llena bien las 2 páginas con texto + 1 elemento visual (lista o imagen).}
\milcpaso{6}{milcMagenta}{\textbf{Paso 6 --- Guarda + dibujo.} Guarda como \texttt{2026-mi-mini-informe.docx} en \emph{Tecnologia/Periodo-3}. En el cuaderno haz un dibujo de las 3 páginas del documento: \emph{página 1 (portada con título grande), página 2 (con encabezado arriba, contenido al medio, pie y \#1 abajo), página 3 (igual con \#2)}.}

\begin{drawbox}{milcMagenta}{Antes de mostrar tu mini-informe}
\checkbox La portada tiene título grande, nombre, grado, fecha. \\[1mm]
\checkbox La portada NO tiene encabezado, pie ni número de página. \\[1mm]
\checkbox Las páginas 2 y 3 tienen encabezado con mi nombre, grado y materia. \\[1mm]
\checkbox Las páginas 2 y 3 tienen pie con datos institucionales. \\[1mm]
\checkbox Las páginas 2 y 3 tienen número de página en esquina inferior derecha. \\[1mm]
\checkbox El documento está guardado con nombre correcto.
\end{drawbox}

\begin{softbox}{milcMagenta}{milcGris}{Plantilla de trabajo}
\textbf{Estructura del informe.} \textit{Página 1 (portada):} título grande + nombre + grado + fecha. \textit{Página 2:} encabezado (nombre·grado·materia) + título de sección + 2 párrafos + pie con datos colegio + número 1. \textit{Página 3:} encabezado igual + cierre del tema + pie + número 2. \textit{Guardado:} 2026-mi-mini-informe.docx.
\end{softbox}

\begin{infoband}{milcMagenta}


\textbf{Lo que va al cuaderno (Actividad 3 · APLICA).} Título: «Actividad 3 --- Mi mini-informe de 2 páginas con marco completo». \textbf{Formato:} dibujo de las 3 páginas del documento etiquetando los 3 elementos en cada una. \textbf{Extensión:} 1 página del cuaderno. \textbf{Sabes que terminaste cuando} cada página del documento tiene los elementos correctos en el lugar correcto.
\end{infoband}

\puente{El producto es ese informe + el boceto en el cuaderno.}

\milcestacion{milcMostaza}{Producto de la sesión}
\textbf{Mini-informe de 2 páginas + portada con marco profesional completo}.
\begin{softbox}{milcMostaza}{milcGris}{Criterios de calidad}
(1) La portada está separada y sin elementos de marco. (2) Las páginas 2 y 3 tienen encabezado con identificación. (3) Las páginas 2 y 3 tienen pie de página con datos institucionales. (4) Las páginas 2 y 3 tienen número de página. (5) El contenido del informe ocupa 2 páginas con texto formateado.
\end{softbox}

\puente{Antes de salir, verificas que los 3 elementos aparecen en las páginas correctas.}

\milcestacion{milcVino}{Evaluación liberadora · cinco dimensiones}

\begin{softbox}{milcVino}{milcGris}{Sentido de la evaluación}
Evaluar no es solo revisar una nota. Es reconocer qué aprendí, qué cambió en mí, qué emociones manejé, cómo me relaciono con la ciudad y la comunidad, y qué le dejaré a quien viene detrás.
\end{softbox}

\textbf{Mi reflexión liberadora en cinco dimensiones.} Completa cada frase iniciada:

\small
\begin{tabularx}{\textwidth}{>{\bfseries\RaggedRight\arraybackslash}p{3.4cm}Y>{\RaggedRight\arraybackslash}p{4.6cm}}
\rowcolor{milcVino}\color{white}Dimensión & \color{white}\bfseries Mi respuesta (frame starter) & \color{white}\bfseries Algo que descubrí\\
1. Personal & Lo que más cambió en mí fue \linead{6.5cm} & \linead{4.2cm}\\
2. Emocional & La emoción más fuerte fue \linead{2.5cm} y la regulé \linead{3cm} & \linead{4.2cm}\\
3. Ciudadana & En democracia esto importa porque \linead{6cm} & \linead{4.2cm}\\
4. Local (Cartago) & En mi barrio sería útil para \linead{6.5cm} & \linead{4.2cm}\\
5. Intergeneracional & A mi abuela/abuelo le diría \linead{6.5cm} & \linead{4.2cm}\\
\end{tabularx}
\normalsize

\begin{infoband}{milcVino}
\textbf{Para la próxima sustentación.} Vuelve a esta tabla en seis meses. Verás que las respuestas cambian — no porque lo aprendido cambie, sino porque tú cambias. La evaluación liberadora es una conversación contigo a través del tiempo.
\end{infoband}


\puente{Cerramos con 3 voces que te ayudan a entender por qué el marco habla antes que el contenido.}

\milcestacion{milcGradeDark}{Triángulo de pensamiento · cierre filosófico}

\begin{softbox}{milcVino}{milcGris}{Dussel · lente del nosotros}
\textit{Nada existe sin marco. Sin marco la pintura es un papel. Sin marco el documento es ruido. El marco da existencia.}

{\footnotesize\itshape\color{milcNegro!70}--- formulación de esta guía, al modo de Enrique Dussel. No es una cita textual.}

Un documento sin marco \textbf{flota sin contexto}: ¿quién lo escribió?, ¿cuándo?, ¿de qué tema?, ¿para qué materia? El marco (encabezado + pie + número) responde esas preguntas \emph{de un vistazo}. Doña Mercedes lo entendía: el cuaderno con nombre, fecha y página era \textbf{un cuaderno con identidad}. Sin esos datos, era un papel cualquiera.

\textbf{¿Mis trabajos tienen marco o flotan sin contexto?}
\end{softbox}
\linead{\linewidth}\\[1mm]
\linead{\linewidth}

\begin{softbox}{milcMostaza}{milcGris}{Estoicismo · Marco Aurelio · cuidado interior}
\textit{La disciplina pequeña hace al trabajo grande visible. Quien marca cada hoja con su nombre se respeta y respeta a quien lo lee.}

{\footnotesize\itshape\color{milcNegro!70}--- formulación de esta guía, al modo de Marco Aurelio. No es una cita textual.}

Poner encabezado y pie toma \textbf{30 segundos por documento}, pero comunica \emph{respeto} y \emph{profesionalismo}. Un profe que recibe 30 trabajos con marco trabaja mejor que uno que recibe 30 sin nombre. \emph{Tu pequeña disciplina les ahorra horas a los demás}. Eso es ética en pequeño.

\textbf{Las pequeñas disciplinas (marco, formato, organización), ¿las cumplo o las dejo para otra vez?}
\end{softbox}
\linead{\linewidth}\\[1mm]
\linead{\linewidth}

\begin{softbox}{milcTurquesa}{milcGris}{Floridi · lente de la infoesfera}
\textit{Un documento con marco es un objeto profesional. Un documento sin marco es un borrador. La diferencia es 30 segundos de configuración.}

{\footnotesize\itshape\color{milcNegro!70}--- formulación de esta guía, al modo de Luciano Floridi. No es una cita textual.}

Hay una diferencia entre \emph{escribir} y \emph{producir un documento profesional}. \textbf{Lo primero es contenido}; \textbf{lo segundo es contenido + marco}. En tu vida adulta, lo que entregas en universidad, trabajo, proyectos, va a juzgarse por las dos cosas. \emph{Hoy aprendiste cómo invertir 30 segundos para subir el nivel}.

\textbf{¿Estoy entregando contenido o estoy entregando documentos profesionales?}
\end{softbox}
\linead{\linewidth}\\[1mm]
\linead{\linewidth}

\begin{infoband}{milcNegro}
\textbf{Nota para el docente.} Las tres formulaciones son del autor de la guía, no citas textuales de los pensadores: recogen su lente para pensar el tema, no sus palabras. Si vas a citarlos en clase, remite a la obra original.
\end{infoband}

\milcestacion{milcVino}{Mi autoevaluación · marca una casilla por fila}

{\small
\setlength{\extrarowheight}{2.2pt}
\begin{tabularx}{\textwidth}{>{\RaggedRight\arraybackslash\bfseries}p{3.0cm}YYY}
\rowcolor{milcVino}\color{white}\bfseries Criterio & \color{white}\bfseries Lo logré & \color{white}\bfseries En proceso & \color{white}\bfseries Necesito apoyo\\[.6mm]
Comprensión del tema & \milcopcion{Explico las ideas con mis palabras.} & \milcopcion{Reconozco las ideas, falta precisión.} & \milcopcion{Necesito repasar lo básico.}\\[.8mm]
\rowcolor{milcGris}Pensamiento computacional & \milcopcion{Usé los pilares al armar mi producto.} & \milcopcion{Usé algunos pilares.} & \milcopcion{Necesito releer la estación 2.}\\[.8mm]
Aplicación y producto & \milcopcion{Mi producto cumple los criterios.} & \milcopcion{Está armado, con ajustes pendientes.} & \milcopcion{Aún está en construcción.}\\[.8mm]
\rowcolor{milcGris}Reflexión 5 dimensiones & \milcopcion{Respondí cada una con honestidad.} & \milcopcion{Respondí algunas con detalle.} & \milcopcion{Mis respuestas son muy generales.}\\[.8mm]
Triángulo de pensamiento & \milcopcion{Conecté las 3 voces con mi trabajo.} & \milcopcion{Conecté al menos una.} & \milcopcion{No las relaciono con mi proceso.}\\[.8mm]
\end{tabularx}}

\vspace{3mm}
\textbf{Hoy entendí que:}\\[1mm]
\linead{\linewidth}\\[1mm]
\linead{\linewidth}

\textbf{Mi compromiso para la próxima sesión es:}\\[1mm]
\linead{\linewidth}\\[1mm]
\linead{\linewidth}

\begin{softbox}{milcMostaza}{milcGris}{Cita de despedida · Marco Aurelio}
\textit{``El obstáculo en el camino se vuelve el camino. Lo que estorba al camino se vuelve camino.''}\\[1mm]
Cada error de hoy, bien observado, se convierte en pieza del camino que pisarás mañana.
\end{softbox}

\small
\begin{tabularx}{\textwidth}{>{\bfseries}p{3.6cm}Y}
\rowcolor{milcGradeDark!12}Nombre completo: & \linead{10cm}\\
Fecha: & \linead{4cm} \quad Curso/grupo: \linead{4cm}\\
Mi firma: & \linead{9cm}\\
\end{tabularx}
\normalsize

\begin{infoband}{milcGradeDark}
\textbf{Para la próxima sesión.} Esta semana, en cualquier documento de más de 1 página que entregues, agrega los 3 elementos de marco. La próxima clase salimos del procesador de texto e iniciamos la \textbf{segunda mitad del periodo}: \emph{Internet}. Vas a entender cómo funciona la red por dentro y cómo viajan los datos cuando navegas. Hoy cerraste el módulo de Word; la próxima abrimos el de Internet.
\end{infoband}

\end{document}
